\begin{proposition}[the real locus]\label{ari:C.2.2}
An antiholomorphic equivalence $X_c\to X_{c'}$ exists exactly when
$c'-(1-\bar c)\in2\mathbb Z$. Consequently $X_c$ admits an
antiholomorphic involution exactly when $q_2\in\mathbb R\setminus\{0\}$.
\end{proposition}
\Uses{the hypotheses, notation and definitions stated here; \cref{fam:src:B.4.2}; \cref{fam:src:A.5.1}; \cite{DuGuoKusakabeXie2026v1}}

\begin{proof}
We check the additional reflection argument, rather than infer a real
structure from the reality of one elliptic invariant. On the period
cover, reflection acts by
\[
 (\tau,\mu,\beta)\longmapsto(-\bar\tau,1-\bar\mu,-\bar\beta).
\]
The theta product and the cusp normalization prove this equality of
continued period functions. For the integral monodromy matrices
$A_1,A_2$ of the retained construction, the lattice comparison is
\[
 H=\begin{pmatrix}1&0&0&0\\0&1&0&0\\-6&0&-1&0\\0&-1&0&-1\end{pmatrix},
 \quad H^2=I,\quad HA_2H=A_2^{-1},\quad
 HA_1H=A_2^{-1}A_1^{-1}A_2,
\]
\[
 \Pi_{-\bar c}(\iota z)H=-\overline{\Pi_c(z)}.
\]
Thus $(z,\zeta)\mapsto(\iota z,-\bar\zeta)$ is an actual torus-family
map. Write $d=(0,0,0,1)^t$, $\varepsilon=(1,2,-4,0)^t$ and
$\varepsilon'=(1,3,-3,0)^t$. Reversing the finite deck generators
changes the target affine vectors to
\[
 v'_1=-\varepsilon+d,\qquad v'_2=\varepsilon'-3d.
\]
The second equality includes the sign of the inverse affine action.
The file \path{Cspecial_reflection.py} checks these matrix identities
over the integers.

Normalize this reflected datum to the retained vectors
$(\varepsilon,-\varepsilon')$. The linear fibre map is $-I_2$, the
parameter shift is $1$, and the lattice map is $-I+E$,
$E=d(1,0,0,0)$. The differences between target and transported affine
vectors are $2d$ and $-4d$. They lie in the images of the respective
norm operators. Explicit integral correcting cocycles are
$\ell_1=(0,0,-2,0)^t$ and $\ell_2=(0,2,0,0)^t$; the norm equations
are $N_1\ell_1=-2d$ and $N_2\ell_2=4d$. Rational solutions of
$(I-A_j)r_j=(v_j-(-I+E)v'_j)/m_j+\ell_j$ can be chosen with first
coordinate zero: $r_1=(0,-2/3,-2/3,0)^t$ and
$r_2=(0,1,1,0)^t$. These give translations holomorphic across both
root discs. Their cusp cocycle is
$\ell_0=-A_0(\ell_1+A_1\ell_2)=(0,0,4,-2)^t$; its coordinates are
integral, so the corresponding shear and inversion preserve the
periodic triangular cusp fan. They therefore extend across its toric
charts. The differences of the local corrections form a cocycle in
the rank-two vector bundle whose invariant line and quotient are
$\mathcal O$ and $\mathcal O(-1)$ on the base sphere. Both have zero
first cohomology, so the local maps glue. Equivalently one splits the
Laurent series of the cocycle into its nonnegative and negative
parts in these two line-bundle frames. This checks the precise
gluing step of Proposition~\ref{fam:src:B.4.2}, including the finite affine translations
and cusp extension. The resulting map has target $X_{1-\bar c}$.

Every other antiholomorphic comparison composed with its inverse is
holomorphic, so \cref{fam:src:A.5.1} gives exactly the congruence in the statement.
A self-map exists iff $\operatorname{Re}c\in1/2+\mathbb Z$, iff
$w=-\bar w$, iff $q_2=iw/216$ is real. Its square is a holomorphic
translation $d_0\in\mathbb C/\mathbb Z$, since the automorphism group
is the $\mathbb C^*$ computed in the credited re-derivation
of Proposition~\ref{fam:src:B.4.1}, Appendix~\ref{app:rederive:fam-src-B.4.1}. The constructed reflection acts on
this additive translation coordinate by conjugation. Commutation
with its square gives $d_0=\bar d_0$ in $\mathbb C/\mathbb Z$, hence
$d_0$ has a real representative. Composing the reflection with
translation by $-d_0/2$ makes its square the identity. Necessity and
sufficiency for an involution follow.
\end{proof}